\section{Implications for Systems Design}
\label{sec:implications}

Each corollary below traces back to a violation of one of the three
well-definedness conditions from Section~\ref{sec:definition}---which is
the point of stating those conditions formally in the first place: once a
defect can be named as a violation of Condition 1, 2, or 3, it stops being
a one-off bug and becomes an instance of a known pattern.

\subsection{Schema Drift as a Domain Mismatch}

\begin{corollary}
\label{cor:schema-drift}
Two components that operate on data with different index sets
$I_1 \neq I_2$, while assuming $I_1 = I_2$, will silently diverge in
behaviour whenever $I_1 \triangle I_2 \neq \emptyset$, regardless of
agreement on $V$.
\end{corollary}

\begin{example}
A producer service serializes records as a function
$D : \{\mathtt{id}, \mathtt{name}, \mathtt{email}\} \to V$. A consumer,
still built against an older schema, deserializes under the assumption
$I' = \{\mathtt{id}, \mathtt{name}\}$. Both sides agree on $V$ and on
every value at every shared index, so nothing looks wrong at the value
level---the defect is purely a mismatch in domain
($\mathtt{email} \in I \setminus I'$), and it stays invisible right up
until some downstream consumer actually needs the \texttt{email} index.
Framing the bug this way says exactly what monitoring ought to check:
not whether the values agree, but whether $I_1 = I_2$ \emph{before} any
value-level comparison is even attempted.
\end{example}

\subsection{Type Confusion as a Codomain Violation}

\begin{corollary}
\label{cor:type-confusion}
A fixed byte-level representation reinterpreted under value sets
$V \neq V'$ for the same indices is not a disagreement about content;
it is the construction of two distinct, both internally well-defined,
functions $D : I \to V$ and $D' : I \to V'$.
\end{corollary}

\begin{example}
Take the same five bytes from Section~\ref{sec:definition}. Under
$V = \text{ASCII codepoints}$, they form the string \texttt{"ARLIZ"};
under $V' = \text{unsigned 8-bit integers}$, the same bytes form the
tuple $(65, 82, 76, 73, 90)$. Neither reading is wrong---$D$ and $D'$
each independently satisfy Conditions 1 through 3 on their own. The
defect only appears when one part of a system builds $D$ under one
codomain and a different part---or, worse, an attacker-controlled input
path---consumes the identical bytes under $V'$ without that change ever
being checked. This is, structurally, the same pattern behind a large
share of deserialization and memory-safety vulnerabilities~\autocite{cwe843}.
\end{example}

\subsection{Unsound Equality from Representational Identity}

\begin{corollary}
\label{cor:equality}
Representational identity (byte-for-byte equality of an encoding) is
necessary but not sufficient for functional equality
(Definition~\ref{def:equality}); systems that conflate the two will
classify distinct collections of data as equal whenever they share a
codomain-independent encoding.
\end{corollary}

\begin{example}
A caching layer that hashes raw bytes to detect duplicate records will
quietly conflate $D$ and $D'$ from the previous example, treating the
string \texttt{"ARLIZ"} and the tuple $(65, 82, 76, 73, 90)$ as the same
cache entry. The corollary points at exactly which piece of structure the
cache's equality check is missing---the codomain $V$---which is also
exactly what the fix has to add: an equality check over data has to take
$V$ into account, not just the bytes.
\end{example}

\subsection{Serialization Ambiguity as Under-Specified Domain or Codomain}

\begin{corollary}
\label{cor:serialization}
A wire format specification that does not fully determine $I$ and $V$ for a
given message admits multiple, individually well-defined, mutually
incompatible deserializations of the same byte stream---a defect that
conformance testing against a single reference implementation cannot detect.
\end{corollary}

\begin{example}
Consider a binary protocol that pins down field positions---fixing
$I$---but leaves a numeric field's signedness unspecified, leaving $V$
under-constrained. A producer that writes the field as a signed integer
and a consumer that reads it as unsigned are both, individually,
implementing the spec correctly. Neither side is buggy; the
specification itself is the defect, and it stops being a defect exactly
when it pins down $I$ and $V$ completely, with nothing left implicit.
\end{example}

\subsection{Summary}

Table~\ref{tab:corollaries} lays out the correspondence between each
structural condition and the failure mode it produces when violated.

\begin{table}[htbp]
\centering
\caption{Structural conditions and their associated failure modes.}
\label{tab:corollaries}
\begin{tabular}{lll}
\hline
\textbf{Violated Condition} & \textbf{Failure Mode} & \textbf{Corollary} \\
\hline
Totality / domain agreement   & Schema drift           & \ref{cor:schema-drift} \\
Codomain consistency          & Type confusion         & \ref{cor:type-confusion} \\
Functional vs.\ rep.\ equality & Unsound equality       & \ref{cor:equality} \\
Domain or codomain under-spec & Serialization ambiguity & \ref{cor:serialization} \\
\hline
\end{tabular}
\end{table}